% ==============================================================
% Adversarial Review Protocol (AR)
% Forschungsprogramm Ontophänomenalismus
% Version 1.0
% Kompiliert mit: pdfLaTeX
% ==============================================================

\documentclass[11pt,a4paper]{article}

% ---------- Grundpakete ----------
\usepackage[utf8]{inputenc}
\usepackage[T1]{fontenc}
\usepackage[german]{babel}
\usepackage{amsmath,amssymb,amsthm}
\usepackage{geometry}
\usepackage{parskip}
\usepackage{enumitem}
\usepackage{booktabs}
\usepackage{xcolor}
\usepackage{hyperref}

% ---------- Layout ----------
\geometry{margin=2.5cm}
\setlist{itemsep=0.2em, topsep=0.2em}
\hypersetup{
    colorlinks=true,
    linkcolor=blue,
    urlcolor=blue,
    pdftitle={Adversarial Review Protocol (AR) -- Version 1.0},
    pdfauthor={Forschungsprogramm Ontoph\"{a}nomenalismus}
}

% ---------- Theorem-Umgebungen ----------
\theoremstyle{definition}
\newtheorem{aufgabe}{Aufgabe}[section]
\newtheorem{regel}{Regel}[section]
\newtheorem{methodennote}{Methodennote}[section]
\newtheorem{remark}{Bemerkung}[section]

% ---------- Titel ----------
\title{
  \textbf{Adversarial Review Protocol (AR)}\\[0.4em]
  {\large Anleitung und Vorlage f\"{u}r externe Reviewer}\\[0.5em]
  {\normalsize Version~1.0}
}
\author{
  \textbf{Forschungsprogramm Ontoph\"{a}nomenalismus}\\
  \small Siehe TRANSPARENCY f\"{u}r methodologische Urheberschaft und KI-Einsatz.
}
\date{Juli 2026}

\begin{document}

\maketitle

\begin{abstract}
Dieses Dokument definiert das Protokoll f\"{u}r externe Adversarial Reviews des
Forschungsprogramms Ontoph\"{a}nomenalismus (OP). Es richtet sich an alle Reviewer ---
menschliche wie k\"{u}nstliche --- und legt Rolle, Vorgehensweise und Output-Format
verbindlich fest.

Ein Adversarial Review ist kein konstruktiver Verbesserungsvorschlag, sondern eine
systematische Suche nach Fehlern, Widerspr\"{u}chen, unbelegten Setzungen und logischen
L\"{u}cken. Das Ziel ist Robustheit: Was einem rigorosen Angriff standhält, ist solide.
\end{abstract}

\tableofcontents
\newpage

%==============================================================
\section{Die Rolle des Adversarial Reviewers}
%==============================================================

\begin{aufgabe}[Rollenverst\"{a}ndnis]
Der Adversarial Reviewer \"{u}bernimmt die Rolle eines informierten, philosophisch und
wissenschaftlich geschulten Kritikers. Er ist wohlwollend gegen\"{u}ber der Intention des
Projekts --- aber kompromisslos gegen\"{u}ber seinen Argumenten.

Seine Aufgabe ist \textbf{nicht}:
\begin{itemize}
  \item Bessere Formulierungen vorschlagen
  \item Das System weiterentwickeln
  \item Zustimmung geben oder loben
\end{itemize}

Seine Aufgabe \textbf{ist}:
\begin{itemize}
  \item Logische Fehler und Widersprüche benennen
  \item Unbelegte Setzungen aufdecken
  \item Kategorienfehler identifizieren
  \item Inkonsistenzen zwischen Dokumenten aufzeigen
  \item Offene Flanken gegen externe Kritik markieren
  \item Festhalten was er f\"{u}r solide h\"{a}lt --- und warum
\end{itemize}
\end{aufgabe}

\begin{methodennote}[Trennung von Review und L\"{o}sung]
Das Review-Dokument benennt Probleme. Es l\"{o}st sie nicht. Die Ausarbeitung von
L\"{o}sungen erfolgt in separaten Diskussionsrunden mit dem Team. Diese Trennung h\"{a}lt
das Review scharf und verhindert, dass Problemerkennung und L\"{o}sungsvorschlag
vermischt werden.
\end{methodennote}

%==============================================================
\section{Vorgehensweise}
%==============================================================

\begin{regel}[Schritt 1: Dokumente lesen]
Der Reviewer liest alle Dokumente des Forschungsprogramms in der folgenden Reihenfolge:

\begin{enumerate}
  \item \textbf{TRANSPARENCY} --- Urheberschaft und Methodik des Projekts
  \item \textbf{MKO} --- Methodologische Konventionen (verbindliche Grundlage)
  \item \textbf{OP} --- Logische Grundlegung (Axiome, Kernarchitektur)
  \item \textbf{OPP} --- Ontoph\"{a}nomenaler Phasenraum (Modul 1)
  \item \textbf{OAD} --- Ontologische Attraktoren-Dynamik (Modul 2)
  \item \textbf{DPR} --- Die Projektive Reduktion (Modul 3)
  \item \textbf{DPD} --- Die Ph\"{a}nomenale Differenz (Modul 4)
  \item \textbf{DFD} --- Das Fragedifferential
  \item \textbf{EG-1} --- Der Explanatory Gap
  \item \textbf{DHP} --- Dekonstruktion des Hard Problems
\end{enumerate}

Die Links zu allen Dokumenten finden sich in der Datei \texttt{Documents.txt} im
Repository.
\end{regel}

\begin{regel}[Schritt 2: MKO als Ma\ss{}stab]
Die Methodologischen Konventionen (MKO) sind der verbindliche Ma\ss{}stab f\"{u}r das
Review. Der Reviewer pr\"{u}ft insbesondere:
\begin{itemize}
  \item Werden Axiome und Folgerungen korrekt im Indikativ formuliert?
  \item Werden Interpretationen im Konjunktiv gehalten?
  \item Sind Analogien explizit als solche gekennzeichnet?
  \item Sind Hypothesen als pr\"{u}fbare Vorschl\"{a}ge ausgewiesen?
  \item Werden emergente Begriffe (Raum, Zeit, Energie) nicht in die ontische Basis
        (OPP, OAD) eingeschmuggelt?
  \item Ist die R/P-Unterscheidung konsequent eingehalten?
\end{itemize}
\end{regel}

\begin{regel}[Schritt 3: Priorisierung der Befunde]
Jeder Befund wird einer von drei Priorit\"{a}tsstufen zugeordnet:

\begin{description}
  \item[\textbf{[K] Kritisch}] Logischer Fehler, echter Widerspruch, kategorialer
        Fehler oder unbelegte Setzung, die das Fundament des Arguments untergr\"{a}bt.
        Muss vor der n\"{a}chsten Reviewrunde adressiert werden.
  \item[\textbf{[B] Bedeutsam}] Argumentative Schw\"{a}che, unvollst\"{a}ndige
        Begr\"{u}ndung oder offene Flanke, die externe Kritik einladen w\"{u}rde.
        Sollte adressiert werden.
  \item[\textbf{[M] Minor}] Terminologische Inkonsistenz, stilistische Schwäche oder
        verbesserungsw\"{u}rdige Formulierung ohne Auswirkung auf das Argument.
\end{description}
\end{regel}

\begin{regel}[Schritt 4: Output erstellen]
Der Reviewer erstellt ein Review-Dokument gem\"{a}\ss{} dem Vorlagen-Format in Sektion~4
dieses Dokuments. Das Ergebnis wird als vollst\"{a}ndiger LaTeX-Code in einem Code-Block
ausgegeben.
\end{regel}

%==============================================================
\section{Verbindliche Review-Kriterien}
%==============================================================

\begin{regel}[Logische Konsistenz]
Folgen die Propositionen und Theoreme aus den Axiomen und Definitionen? Gibt es
versteckte Pr\"{a}missen? Gibt es Zirkelschl\"{u}sse?
\end{regel}

\begin{regel}[Terminologische Konsistenz]
Werden dieselben Begriffe in allen Dokumenten einheitlich verwendet? Gibt es
Bedeutungsverschiebungen zwischen Modulen?
\end{regel}

\begin{regel}[Kontaminationskontrolle]
Werden emergente Begriffe (Raum, Zeit, Energie, Abstand, Bewegung, Symbiose) in
OPP oder OAD als konstituierende Eigenschaften verwendet? Solche Stellen sind
als [K] kritisch zu markieren.
\end{regel}

\begin{regel}[Modulare Koherenz]
Sind die Schnittstellen zwischen den Modulen konsistent? Liefert OAD wirklich was
DPR voraussetzt? Liefert DPR wirklich was DPD und DFD voraussetzen?
\end{regel}

\begin{regel}[Epistemische Kennzeichnung]
Sind alle Aussagen korrekt nach MKO-Ebene gekennzeichnet? Gibt es Stellen wo
Indikativ verwendet wird, aber Konjunktiv erforderlich w\"{a}re?
\end{regel}

\begin{regel}[Externe Angreifbarkeit]
Wo w\"{u}rde ein analytischer Philosoph, ein Physiker oder ein Kognitionswissenschaftler
sofort einhaken? Diese Stellen sollen explizit markiert werden --- auch wenn das
System intern konsistent ist.
\end{regel}

\begin{regel}[Positive Bewertung]
Der Reviewer benennt explizit, was er f\"{u}r argumentativ solide h\"{a}lt. Ein Review
ohne positive Bewertungen ist unvollst\"{a}ndig. Die Begr\"{u}ndung muss angegeben werden:
Nicht \glqq{}das ist gut\grqq{}, sondern \glqq{}das ist gut, weil\ldots\grqq{}.
\end{regel}

%==============================================================
\section{Output-Vorlage f\"{u}r das Review-Dokument}
%==============================================================

Der Reviewer erstellt ein Dokument nach dem folgenden Schema. Der Name des Reviewers
und das Datum werden im Titel eingetragen.

\begin{methodennote}[Verwendung der Vorlage]
Die folgende Vorlage ist verbindlich. Abweichungen im Format sind nur erlaubt wenn
sie begr\"{u}ndet werden. Der gesamte Output wird als vollst\"{a}ndiger, kompilierbarer
LaTeX-Code ausgegeben.
\end{methodennote}

\newpage
\noindent\textbf{Vorlage f\"{u}r das Review-Dokument:}

\begin{verbatim}
% ==============================================================
% Adversarial Review: Ontophänomenalismus (OP)
% Reviewer: [Name des Reviewers]
% Datum: [Datum]
% ==============================================================

\documentclass[11pt,a4paper]{article}
\usepackage[utf8]{inputenc}
\usepackage[T1]{fontenc}
\usepackage[german]{babel}
\usepackage{geometry}
\usepackage{parskip}
\usepackage{enumitem}
\usepackage{booktabs}
\usepackage{hyperref}
\geometry{margin=2.5cm}

\title{
  \textbf{Adversarial Review}\\
  \textbf{Forschungsprogramm Ontophänomenalismus}\\[0.4em]
  {\normalsize Reviewer: [Name] --- [Datum]}
}
\author{}
\date{}

\begin{document}
\maketitle

%--------------------------------------------------------------
\section*{Zusammenfassung}
%--------------------------------------------------------------
% 3-5 Sätze: Was ist der Gesamteindruck? Was sind die
% gravierendsten Probleme? Was ist besonders solide?

%--------------------------------------------------------------
\section*{Kritische Befunde [K]}
%--------------------------------------------------------------
% Logische Fehler, echte Widersprüche, kategoriale Fehler,
% unbelegte Setzungen die das Fundament untergraben.
% Format pro Befund:

% \subsection*{[K-01] Titel des Befunds}
% \textbf{Dokument:} [Kürzel und Version]\\
% \textbf{Stelle:} [Sektion oder Satz]\\
% \textbf{Problem:} [Präzise Beschreibung]\\
% \textbf{Warum kritisch:} [Begründung]

%--------------------------------------------------------------
\section*{Bedeutsame Befunde [B]}
%--------------------------------------------------------------
% Argumentative Schwächen, offene Flanken, unvollständige
% Begründungen.
% Format wie [K], aber mit [B-01] etc.

%--------------------------------------------------------------
\section*{Minore Befunde [M]}
%--------------------------------------------------------------
% Terminologische Inkonsistenzen, Formulierungsschwächen.
% Format wie [K], aber mit [M-01] etc.

%--------------------------------------------------------------
\section*{Positive Bewertungen}
%--------------------------------------------------------------
% Was ist argumentativ solide --- und warum?
% Format pro Bewertung:

% \subsection*{[P-01] Titel}
% \textbf{Dokument:} [Kürzel]\\
% \textbf{Stelle:} [Sektion]\\
% \textbf{Stärke:} [Begründung warum solide]

%--------------------------------------------------------------
\section*{Offene Fragen an das Team}
%--------------------------------------------------------------
% Fragen die nicht als Fehler klassifiziert werden können,
% aber einer Klärung bedürfen.

%--------------------------------------------------------------
\section*{Priorisierte Empfehlungen}
%--------------------------------------------------------------
% Welche drei Befunde sollten zuerst adressiert werden?
% Kurze Begründung der Reihenfolge.

\end{document}
\end{verbatim}

%==============================================================
\section{Hinweise f\"{u}r den Reviewer}
%==============================================================

\begin{remark}[Zum Umgang mit dem transzendentalen Zirkel]
Das Forschungsprogramm selbst benennt den transzendentalen Zirkel als unvermeidlich
und arbeitet mit einer expliziten R/P-Unterscheidung. Kritik an der Zirkularit\"{a}t
allein gen\"{u}gt nicht --- der Reviewer muss zeigen, dass der OP den Zirkel
\emph{schlechter} kontrolliert als rivalisierende Systeme.
\end{remark}

\begin{remark}[Zum Umgang mit der PST-Blackbox]
Die Projektive Strukturtheorie (PST) ist bewusst als offenes Forschungsprogramm
ausgewiesen. Die Kritik \glqq{}die PST existiert noch nicht\grqq{} ist kein
g\"{u}ltiger Befund --- das wei\ss{} das Team. Ein g\"{u}ltiger Befund w\"{a}re:
\glqq{}Die Anforderungen an die PST sind zu unscharf um pr\"{u}fbar zu sein\grqq{}
oder \glqq{}Dokument X setzt PST-Ergebnisse voraus, die noch nicht formuliert sind.\grqq{}
\end{remark}

\begin{remark}[Zur Sprache des Reviews]
Das Review-Dokument selbst muss nicht MKO-konform sein. Der Reviewer schreibt in
seiner eigenen Sprache. Klarheit und Pr\"{a}zision sind wichtiger als Einhaltung
der OP-Konventionen.
\end{remark}

\begin{remark}[Zum Umfang]
Qualit\"{a}t vor Quantit\"{a}t. Drei pr\"{a}zise [K]-Befunde sind wertvoller als
zwanzig oberfl\"{a}chliche [M]-Befunde. Der Reviewer soll tief bohren, nicht breit
str\"{a}uen.
\end{remark}

\end{document}

